\scanpage{66}
\mspage{7}
% archive foliation: 68
coïncident. Cela signifie donc que
$X$ est déduit de $X^*$ par ``bouchage
de trous \unsure{par disques} canoniques'' (utilisant
les structures toriques sur les
comp. conn. de $\delta X^*$). \deleted{La question}\inserted{Disons alors}
qu'un système de disques $(D_i)_{i\in I}$ dans $X$
\deleted{\unsure{qui n}} autour des $a_i$ est ``admissible''.

\uline{\unsure{Questions|Question :}} Comment cette condition
d'admissibilité se reflète-t-elle en
termes de la géométrie interne de $X$ ?
Par exemple, \deleted{quand} quand $X$ est muni
d'une métrique de \unsure{Bergmann|Bergman} canonique
(cas où $X\not\simeq\mathbb P^1_{\mathbb C}$), quelle est la
relation de cette condition avec celle
que les \deleted{$\unsure{\delta}$}$D_i$ soient des ``disques
géodésiques'' (\inserted{définis} par la condition que
\[
\operatorname{\text{dist. géod}}(x,a_i)\leq r_i
\]
) ?

\margin{supposons hyperbolique i.e. $(g,\nu)\neq(0,1),(0,2)$}
Fixons une \deleted{$X_0$} surface top. \inserted{orientée}
à bord de référence $X_0$, de type $g,\nu$
